\section{The model M1 and its rationale}
\label{sec:model}

\begin{definition}[Model M1]\label{def:m1}
For an environment $u=(y,t,e)$, write $c(u)=(y,t)\in\Cset$ and $r(u)=(t,e)\in\Rset$ for the two cells it
determines. The latent vector for an image from environment $u$ is $Z^u=(Z_c^u,Z_s^u)$ with
$Z_c^u\in\R^{n_c}$ (`content'), $Z_s^u\in\R^{n_s}$ (`style'), $n=n_c+n_s$, and the image is
$X^u\in\R^{d}$. Note that Ng et al. assume $p^{(u)}_Z$ is third-order differentiable and 
\emph{has full support on $\R^n$}.The data are generated as follows.
\begin{align}
X^u &= g(Z_c^u,Z_s^u), \label{eq:mixing}\\
Z_{c,i}^u &= f_i^{c(u)}\bigl(\PA(Z_{c,i}^u;\Gc)\bigr) + \varepsilon_{c,i}^u,
\qquad \varepsilon_{c,i}^u\sim\Nn\bigl(\mu_{c,i}^{c(u)},(\sigma_{c,i}^{c(u)})^2\bigr),
\quad i=1,\dots,n_c, \label{eq:content}\\
Z_{s,j}^u &= \varepsilon_{s,j}^u,
\qquad \varepsilon_{s,j}^u\sim\Nn\bigl(\mu_{j}^{r(u)},(\sigma_{s,j}^{r(u)})^2\bigr),
\quad j=1,\dots,n_s, \label{eq:style}
\end{align}
Here $\Gc$ is an unknown DAG over the content latents, $g$ is an unknown mixing function, and the noise
variables are mutually independent (within an image, and across images). Neither $\Gc$ nor $g$ carries
an index: they are the same for every environment. Everything that changes with $u$ -- the mechanisms
$f_i^{c(u)}$, $\mu_{c,i}^{c(u)}$, $\sigma_{c,i}^{c(u)}$ and the style parameters $\mu_j^{r(u)}$, $\sigma_{s,j}^{r(u)}$ --
changes only through $c(u)$ or $r(u)$, never through $u$ as a whole; this is the content of
\labelcref{item:P2,item:P3} below. In addition:
\begin{enumerate}[label=(P\arabic*),leftmargin=2.4em]
\item \label{item:P1} \emph{Invariant mixing.} $g$ is a $\mathrm{C}^2$ diffeomorphism onto its image and is the same for all
      species, times of day and locations.
\item \label{item:P2} \emph{Invariant content mechanisms.} The DAG $\Gc$ and the content mechanisms
      $(f_i^{(c)},\mu_{c,i}^{(c)},\sigma_{c,i}^{(c)})_{c\in\Cset}$ depend on the environment $u$ only through $c(u)$, not
      on the location: for environments $u=(y,t,e)$ and $u'=(y,t,e')$ with the same $(y,t)$, $Z_c^u$ and
      $Z_c^{u'}$ have the same law.
\item \label{item:P3} \emph{Exogenous style.} The style latents have no parents and no children among the latents.
      Their distribution depends on the environment $u$ only through $r(u)=(t,e)$, and they are
      independent of the species $Y$ and of $Z_c^u$ given $(t,e)$.
\item \label{item:P4} \emph{Location-specific label prior.} $C\mid E=e\sim p^{e}(C)$ is allowed to differ across
      locations; this is a fact about the observed variables $(Y,T,E)$, external to the structural
      equations \eqref{eq:mixing}--\eqref{eq:style}, and is used only for prediction
      (\cref{sec:use}).
\end{enumerate}
The latent DAG over $Z^u$ is $\Gc$ together with $n_s$ isolated nodes (the style latents), the same DAG
for every $u$. The time of day $T$ enters both blocks: it is part of $c(u)$, which indexes the content
mechanisms, and part of $r(u)$, which indexes the style distributions.
\end{definition}

\Cref{fig:m1} shows the model.

\begin{figure}[H]
\centering
\begin{tikzpicture}[
  >=Stealth,
  node distance=8mm and 12mm,
  latnode/.style={
    circle,draw,minimum size=8mm,inner sep=1pt
  },
  obsnode/.style={
    circle,draw,fill=gray!25,minimum size=8mm,inner sep=1pt
  },
  plate/.style={
    draw,rounded corners,inner sep=4mm
  },
  platelabel/.style={font=\small,anchor=north west,inner sep=2pt},
  paramplate/.style={
    draw,
    rounded corners,
    inner xsep=2mm,
    inner ysep=1mm,
    align=left
  }
]

% ---- observed indices, inside the sample plate ----
\node[obsnode] (Y) at (-1.4,2.6) {$Y$};
\node[obsnode] (T) at (Y |- 0,0.6) {$T$};
\node[obsnode] (E) at (Y |- 0,-1.6) {$e$};

% ---- content latents (one sample's worth), with the shared DAG structure drawn once ----
\node[latnode] (c1) at (1.8,2.8)     {$Z_{c,1}^u$};
\node[latnode] (c2) at (4.2,0 |- c1) {$Z_{c,2}^u$};
\node[latnode] (c3) at (c1 |- 0,1.0) {$Z_{c,3}^u$};
\node[latnode] (c4) at (c2 |- c3)    {$Z_{c,4}^u$};

\draw[->] (c1) -- (c2);
\draw[->] (c1) -- (c3);
\draw[->] (c2) -- (c4);
\draw[->] (c3) -- (c4);

\begin{scope}[on background layer]
  \node[
    fit=(c1)(c2)(c3)(c4),
    draw,dashed,rounded corners,
    inner sep=5mm
  ] (C) {};
\end{scope}

\node[
  font=\scriptsize,
  below=1mm of C.south,
  xshift=-2mm
] {structure $\Gc$ (shared, no plate)};

\node[latnode] (s1) at (3.0,-1.6) {$Z_{s}^u$};
\node[obsnode] (X)  at (6.6,0.6)  {$X^u$};

\draw[->]
  (C.east) -- node[above,sloped,font=\small]{$g$} (X);

\draw[->]
  (s1) -- node[below,sloped,font=\small]{$g$} (X);

% ---- the sample plate: one copy per image / environment u=(y,t,e) ----
\begin{scope}[on background layer]
  \node[
    plate,
    fit=(Y)(T)(E)(C)(s1)(X),
    inner sep=6mm
  ] (plateU) {};
\end{scope}

\node[platelabel] at (plateU.north west) {$u=(y,t,e)$};

% ---- content-mechanism parameter plate: parameters indexed by the content cell c, no sample index ----
\node[
  paramplate
] (plateC) at (3.65,5.9) {%
  \begin{tabular}{@{}l@{}}
    $c\in\Cset_{\mathrm{occ}}$\\[0.1mm]
    $f_i^{(c)},\ \mu_{c,i}^{(c)},\ \sigma_{c,i}^{(c)}$\footnotemark\\[0.1mm]
    $i=1,\dots,n_c$
  \end{tabular}%
};

\draw[->]
  (plateC.south) --
  node[right,font=\small,pos=.3]{$c(u)=(y,t)$}
  (C.north);

% ---- style plate: parameters indexed by the style cell r, no sample index ----
\node[
  paramplate
] (plateR) at (3.65,-4.3) {%
  \begin{tabular}{@{}l@{}}
    $r\in\Rset_{\mathrm{occ}}$\\[0.1mm]
    $\mu_j^{(r)},\ \sigma_{s,j}^{(r)}$\\[0.1mm]
    $j=1,\dots,n_s$
  \end{tabular}%
};

\draw[->]
  (plateR.north) --
  node[right,font=\small,pos=.3]{$r(u)=(t,e)$}
  (s1.south);

% ---- P4: location-specific label prior, external to the SEM ----
\draw[->,dashed]
  (E) to[bend right=25]
  node[below,sloped,font=\scriptsize]
    {$p^e(C)$, prediction only (P4)}
  (Y);

\draw[->,dashed]
  (E) to[bend left=25]
  (T);

\end{tikzpicture}
\caption{\textbf{The M1 model, in plate notation.} The large box is the sample plate: one copy of
$(Y,T,e,Z_c^u,Z_s^u,X^u)$ per image, replicated over $u=(y,t,e)$. The two small boxes above and below
are parameter plates, replicated over the content cell $c=(y,t)$ and the style cell $r=(t,e)$
respectively; nothing inside them carries a sample index. The arrow from a parameter plate into the
sample plate is selected by that \textbf{same} image's own $(y,t)$ or $(t,e)$, as labelled (note that
$t$ is \textbf{shared}). Objects with no plate around them -- the DAG structure $\Gc$ and the mixing
$g$ -- are shared by every sample and every cell.
The dashed arrows are not part of the structural equations \eqref{eq:mixing}--\eqref{eq:style}: they
record the location-specific label prior of \labelcref{item:P4}, used only at prediction time
(\cref{sec:use}). $Y$ is unobserved and $T$ can be either observed or unobserved. The four content latents and their edges are only an illustration: the number of
content latents and their DAG are unknown.}
\label{fig:m1}
\end{figure}
\footnotetext{In an implementation, $\mu_{c,i}^{(c)}$ can be absorbed into $f_i^{(c)}$}.

\paragraph{Rationale.}
\begin{itemize}[leftmargin=1.5em]
\item \textbf{Invariance is placed where it is plausible.} Species and time of day produce animal
      appearance through a mechanism that should not depend on the camera site, so $p(z_c\mid c)$ and
      $g$ are shared across locations---\labelcref{item:P1},\labelcref{item:P2}.
\item \textbf{Locations differ by more than the label prior.} A pure shift of $p^e(C)$ (label shift)
      cannot account for differences of geography, camera and orientation. These are modelled as
      exogenous variables whose distribution depends on the location---\labelcref{item:P3}.
\item \textbf{Style depends on the time of day as well.} A specific camera and its position at a
      location may be affected in a time-of-day-specific way. The style distribution is therefore
      indexed by $(t,e)$ and not by $e$ alone. The time of day also affects the content,
      as in \labelcref{item:P2}, so it enters both
      blocks. Since $Z_s^u$ is independent of the species and of $Z_c^u$ given $(t,e)$, it carries no
      information about the species beyond what the observed $t$ already provides.
\item \textbf{Discreteness is kept in an observed variable.} $C$ is observed, so all \emph{latent}
      variables can be continuous, which is the setting of the continuous theory of
      \cite{ng2025}. No discrete-latent theory is needed.
\item \textbf{The content latents are not assumed independent.} Their DAG $\Gc$ is unknown and
      possibly non-empty; finding it is part of the problem.
\item \textbf{The number of latents is a modelling choice.} Because $C$ only indexes mechanisms,
      $n_c$ and $n_s$ are not tied to $|\Cset|$; they are hyperparameters, bounded above by the
      requirements derived in \cref{sec:reqs}.
\item \textbf{Relation to earlier frameworks.} The split into an invariant block and a changing
      block follows the partial-disentanglement idea of \cite{kong2022}; using an observed auxiliary
      variable to index changing distributions is the mechanism of \cite{khemakhem2020}. This also maps
      naturally to the application with the auxiliary variable incorporating distributions changing
      observables.
\end{itemize}
